\documentclass{article}
\usepackage{amsmath, amssymb, amsthm}

\title{IMC 2025 \\ First Day, July 30, 2025 \\ Solutions}
\date{}

\begin{document}
\maketitle

\section*{Problem 1}
Let $P \in \mathbb{R}[x]$ be a polynomial with real coefficients, and suppose $\deg(P) \geq 2$. For every $x \in \mathbb{R}$, let $\ell_x \subset \mathbb{R}^2$ denote the line tangent to the graph of $P$ at the point $(x, P(x))$.
\begin{itemize}
\item[(a)] Suppose that the degree of $P$ is odd. Show that $\bigcup_{x \in \mathbb{R}} \ell_x = \mathbb{R}^2$.
\item[(b)] Does there exist a polynomial of even degree for which the above equality still holds?
\end{itemize}
(proposed by Mike Daas, Max Planck Institute for Mathematics, Bonn)

\paragraph{Solution.}

(a) Suppose that the degree of $P$ is odd and let $(a,b) \in \mathbb{R}^2$ be arbitrary. Given $r \in \mathbb{R}$, the equation for $\ell_r$ is given by
\[
\ell_r = \{(x,y) \in \mathbb{R}^2 \mid y = P'(r)(x-r) + P(r)\}.
\]
For this line to pass through the point $(a,b)$ it is therefore necessary and sufficient that
\[
b = a P'(r) + P(r) - r P'(r).
\]
This is a polynomial equation in $r$, which always has a real solution as soon as we can show that the degree is odd. Indeed, if $P(r) = c r^n + \ldots$ describes the leading term, then the right hand side of the above equation has leading term $(c - nc) r^n$. Since $c \neq 0$ and we assumed that $n \geq 2$, we must have $c - nc \neq 0$. The right hand side therefore has the same degree as $P$, completing the proof.

(b) If the degree of $P$ is even, then this can never be true, because the degree of $a P'(r) + P(r) - r P'(r)$ is now even by the same argument and therefore it has a global minimum (or maximum) over the reals, below (or above) which no value of $b$ will yield a real solution for $r$.

\section*{Problem 2}
Let $f : \mathbb{R} \to \mathbb{R}$ be a twice continuously differentiable function, and suppose that $\int_{-1}^{1} f(x)\,dx = 0$ and $f(1) = f(-1) = 1$. Prove that
\[
\int_{-1}^{1} (f''(x))^2 \, dx \geq 15,
\]
and find all such functions for which equality holds.

(proposed by Alberto Cagnetta, Universit\`a degli Studi di Udine, Italy)

\paragraph{Solution.}
If $g$ is an arbitrary twice continuously differentiable function on $[-1,1]$, then by applying the Cauchy--Schwarz inequality for $f''$ and $g$ and integrating by parts twice we get
\begin{align*}
\sqrt{\int_{-1}^{1} (f'')^2 \cdot \int_{-1}^{1} g^2} &\geq \int_{-1}^{1} f'' g = \bigl(f'(1)g(1) - f'(-1)g(-1)\bigr) - \int_{-1}^{1} f' g' \\
&= \bigl(f'(1)g(1) - f'(-1)g(-1)\bigr) - \bigl(f(1)g'(1) - f(-1)g'(-1)\bigr) + \int_{-1}^{1} f g'' \\
&= \bigl(f'(1)g(1) - f'(-1)g(-1) - g'(1) + g'(-1)\bigr) + \int_{-1}^{1} f g''. \tag{1}
\end{align*}
In order to get rid of the terms $f'(1)g(1)$ and $f'(-1)g(-1)$ we choose $g$ such that $g(1) = g(-1) = 0$. Moreover, if $g''$ is constant, so $g$ is an at most quadratic polynomial, then $\int_{-1}^{1} f g'' = g'' \int_{-1}^{1} f = 0$.
Hence, it is reasonable to apply (1) with $g(x) = (1+x)(1-x) = 1 - x^2$.

With this choice,
\[
g(1) = g(-1) = 0, \quad g'(1) = -2, \quad g'(-1) = 2, \quad g'' \equiv -2,
\]
and
\[
\int_{-1}^{1} g^2 = \int_{-1}^{1} (1 - 2x^2 + x^4)\, dx = \frac{16}{15},
\]
so we get
\[
\sqrt{\int_{-1}^{1} (f'')^2 \cdot \frac{16}{15}} \geq 0 - 0 - g'(1) + g'(-1) + 0 = 4,
\]
\[
\int_{-1}^{1} (f'')^2 \geq 15.
\]

Equality in Cauchy--Schwarz in (1) holds only if there exists a real $\lambda$ such that $f''(x) = \lambda g(x)$ almost everywhere in $[-1,1]$; by continuity of $f''$ and $g$ we have $f'' = \lambda g$ everywhere on the interval $[-1,1]$. Hence $f(x)$ must have the form
\[
f(x) = \lambda\left(\frac{x^2}{2} - \frac{x^4}{12}\right) + ax + b, \quad \text{with } \lambda, a, b \in \mathbb{R}.
\]
From $\int_{-1}^{1} f(x)\,dx = 0$, we get $0 = \frac{3\lambda}{10} + 2b$, hence $b = -\frac{3}{20}\lambda$. Moreover, the condition $f(1) = f(-1) = 1$ implies
\[
\lambda\left(\frac{1}{2} - \frac{1}{12}\right) + a - \frac{3}{20}\lambda = f(1) = 1 = f(-1) = \lambda\left(\frac{1}{2} - \frac{1}{12}\right) - a - \frac{3}{20}\lambda
\]
hence $a = 0$, and $\lambda = \frac{15}{4}$.

In conclusion, equality holds if and only if
\[
f(x) = \frac{15}{4}\cdot\left(\frac{x^2}{2} - \frac{x^4}{12} - \frac{3}{20}\right) = \frac{-5x^4 + 30x^2 - 9}{16} \quad \text{for all } x \in [-1,1].
\]

\section*{Problem 3}
Denote by $\mathcal{S}$ the set of all real symmetric $2025 \times 2025$ matrices of rank 1 whose entries take values $-1$ or $+1$. Let $A, B \in \mathcal{S}$ be matrices chosen independently uniformly at random. Find the probability that $A$ and $B$ commute, i.e.\ $AB = BA$.

(proposed by Marian Panţiruc, ``Gheorghe Asachi'' Technical University of Iaşi, Romania)

\paragraph{Solution.}
Let $n = 2025$. First, we give a characterisation of matrices in $\mathcal{S}$.

Suppose that $A = (a_{ij})_{i,j=1}^n \in \mathcal{S}$. Since $\mathrm{rk}\, A = 1$, for every $1 < i,j \leq n$, we have
\[
\det\begin{pmatrix} a_{11} & a_{1j} \\ a_{i1} & a_{ij} \end{pmatrix} = a_{11}a_{ij} - a_{i1}a_{1j} = a_{11}a_{ij} - a_{i1}a_{j1} = 0.
\]
If $a_{11} = 1$ then this means $a_{ij} = a_{i1} a_{j1}$. In this case, let $u = (1, a_{21}, \ldots, a_{n1})^{\top}$; then we have $A = (a_{i1}a_{j1}) = u u^{\top}$. Otherwise, if $a_{11} = -1$, we have $a_{ij} = -a_{i1}a_{j1}$. In that case, let $u = (1, -a_{21}, \ldots, -a_{n1})^{\top}$; then $A = -(a_{i1} a_{j1}) = -u u^{\top}$.

Hence, all matrices in $\mathcal{S}$ can uniquely be written as $\pm u u^{\top}$ with a vector $u = (1, u_2, \ldots, u_n)^{\top}$ such that $u_2, \ldots, u_n \in \{\pm 1\}$. (Note that $\mathrm{rk}(u u^{\top}) = 1$ is satisfied.) In particular, we have $|\mathcal{S}| = 2^n$, because the sign and the coordinates $u_2, \ldots, u_n$ can be chosen independently.

Now, if $A = \pm u u^{\top}$ and $B = \pm v v^{\top}$ are elements of $\mathcal{S}$, then
\[
AB = \pm (u u^{\top})(v v^{\top}) = \pm u(u^{\top} v) v^{\top} = \pm \langle u, v \rangle \cdot (u_i v_j)_{i,j=1}^n
\]
and similarly
\[
BA = \pm \langle u, v \rangle \cdot (v_i u_j)_{i,j=1}^n.
\]
Since $n = 2025$ is odd, it follows that $\langle u, v \rangle \neq 0$. The first columns of the matrices $u v^{\top}$ and $v u^{\top}$ are $u$ and $v$, respectively. Hence, $AB = BA$ if and only if $u = v$; in other words, if $A = \pm B$.

For each $A \in \mathcal{S}$, there are precisely two suitable matrices $B \in \mathcal{S}$, so the probability that $A, B$ commute is
\[
\frac{2}{|\mathcal{S}|} = \frac{1}{2^{n-1}}.
\]

\section*{Problem 4}
Let $a$ be an even positive integer. Find all real numbers $x$ such that
\begin{equation}
\left\lfloor \sqrt[a]{b^a + x} \cdot b^{a-1} \right\rfloor = b^a + \lfloor x/a \rfloor
\end{equation}
holds for every positive integer $b$.

(proposed by Yagub Aliyev, ADA University, Baku, Azerbaijan)

\paragraph{Solution.}
We will show that if $a = 2$ then we must have $x \in [-1, 2) \cup [3, 4)$, otherwise $x \in [-1, a)$.

Suppose that $\lfloor x/a \rfloor = m$. Then $m \leq x/a < m+1$, and
\begin{equation}
am \leq x < a(m+1). \tag{2}
\end{equation}
Let $b = 1$. Then
\begin{equation}
\left\lfloor \sqrt[a]{1 + x} \right\rfloor = 1 + \lfloor x/a \rfloor. \tag{3}
\end{equation}
From (3) it follows that $\left\lfloor \sqrt[a]{1+x} \right\rfloor = 1 + m$, or $1 + m \leq \sqrt[a]{1+x} < 2 + m$, or
\begin{equation}
(1+m)^a - 1 \leq x < (2+m)^a - 1, \tag{4}
\end{equation}
where we have used the obvious fact that $m \geq -1$. Indeed, the number $\sqrt[a]{1+x}$ and therefore the number $\lfloor \sqrt[a]{1+x} \rfloor = 1 + m$ is not negative. By Bernoulli's inequality, the following two inequalities --- which compare the inequalities (2) and (4) --- hold:
\begin{align}
am &\leq (1+m)^a - 1, \tag{5}\\
a(m+1) &\leq (1 + (m+1))^a - 1, \tag{6}
\end{align}
where in (5), equality holds if and only if $m = 0$, and in (6), equality holds if and only if $m = -1$.
From (2), (4)--(6), it follows that
\begin{equation}
(m+1)^a - 1 \leq x < a(m+1). \tag{7}
\end{equation}
From (7), it follows that $(m+1)^a - 1 < a(m+1)$. Therefore
\begin{equation}
(m+1)^a \leq a(m+1). \tag{8}
\end{equation}
From (8), it follows that $m+1 \leq a^{\frac{1}{a-1}}$. If $a > 2$ then $m + 1 < a^{\frac{1}{a-1}} < 2$, because $2^{a-1} > a$ for $a > 2$ (one can prove this by mathematical induction) and $1 < a^{\frac{1}{a-1}} < 2$ is not an integer. Therefore, if $a > 2$ then $m = 0$ or $m = -1$. If $a = 2$ then $m = -1$, $m = 0$ or $m = 1$. From (7), it follows that equality (3) holds true only for the values $-1 \leq x < 0$ ($m = -1$), $0 \leq x < a$ ($m = 0$) if $a > 2$, and $-1 \leq x < 0$ ($m = -1$), $0 \leq x < 2$ ($m = 0$) and $3 \leq x < 4$ ($m = 1$) if $a = 2$.

We will now prove that for these values of $x$, equality (1) is true for all positive integers $b$.
From (7), it follows that $b^a + (m+1)^a - 1 \leq b^a + x < b^a + a(m+1)$ and
\begin{equation}
1 + \frac{(m+1)^a - 1}{b^a} \leq 1 + \frac{x}{b^a} < 1 + \frac{a(m+1)}{b^a}. \tag{9}
\end{equation}
By Bernoulli's inequality
\begin{equation}
1 + \frac{a(m+1)}{b^a} \leq \left(1 + \frac{m+1}{b^a}\right)^a, \tag{10}
\end{equation}
where equality occurs if and only if $m = -1$. It is easy to check that for $m = 1$ (if $a = 2$), $m = 0$ and $m = -1$, the following inequality holds true:
\begin{equation}
\left(1 + \frac{m}{b^a}\right)^a \leq 1 + \frac{(m+1)^a - 1}{b^a}. \tag{11}
\end{equation}
From (9)--(11), it follows that
\begin{equation}
\left(1 + \frac{m}{b^a}\right)^a \leq 1 + \frac{x}{b^a} < \left(1 + \frac{m+1}{b^a}\right)^a. \tag{12}
\end{equation}
From (12), it follows that
\begin{align*}
b^a + m &\leq \sqrt[a]{b^a + x} \cdot b^{a-1} < b^a + m + 1,\\
b^a + \lfloor x/a \rfloor &\leq \sqrt[a]{b^a + x} \cdot b^{a-1} < b^a + \lfloor x/a \rfloor + 1.
\end{align*}
Consequently, equality (1) holds.

\section*{Problem 5}
For a positive integer $n$, let $[n] = \{1, 2, \ldots, n\}$. Denote by $S_n$ the set of all bijections from $[n]$ to $[n]$, and let $T_n$ be the set of all maps from $[n]$ to $[n]$. Define the \emph{order} $\mathrm{ord}(\tau)$ of a map $\tau \in T_n$ as the number of distinct maps in the set $\{\tau, \tau \circ \tau, \tau \circ \tau \circ \tau, \ldots\}$ where $\circ$ denotes composition. Finally, let
\[
f(n) = \max_{\tau \in S_n} \mathrm{ord}(\tau) \quad \text{and} \quad g(n) = \max_{\tau \in T_n} \mathrm{ord}(\tau).
\]
Prove that $g(n) < f(n) + n^{0.501}$ for sufficiently large $n$.

(proposed by Fedor Petrov, St Petersburg State University)

\paragraph{Solution.}
For every $\tau \in T_n$ we need to prove that $\mathrm{ord}(\tau) \leq f(n) + n^{0.501}$ (if $n$ is large enough). Denote by $C(\tau)$ the set of elements $x \in [n]$ for which $\tau^k(x) = x$ for some $k > 0$. It is immediate that $C(\tau)$ is a $\tau$-invariant set; let $\tau_c = \tau|_{C(\tau)}$, that is a permutation on $C(\tau)$.

Let $N$ be the order of this permutation $\tau_c$, clearly $N \leq g(n)$. Consider an arbitrary element $x \in [n] \setminus C(\tau)$. The sequence $x, \tau(x), \tau^2(x), \ldots$ is eventually periodic, but not from the beginning, because $x \notin C(\tau)$.

Let $h(x) > 0$ be the minimal number for which $\tau^{h(x)}(x)$ is in the period; equivalently, this is the minimal number for which $\tau^{h(x)}(x) \in C(\tau)$. Let $R = \max_{x \in [n] \setminus C(\tau)} h(x)$. Note that $\tau^R(x) = \tau^{R+N}(x)$ for all $x \in [n]$, since $\tau^R(x) \in C(\tau)$ for all $x \in [n]$. Therefore, $\mathrm{ord}(\tau) \leq N + R$. Thus, if $R < n^{0.501}$, we are done.

Now assume that $R \geq n^{0.501}$, that is, $h(x) \geq n^{0.501}$ for some $x \notin C(\tau)$. It yields $|C(\tau)| \leq n - n^{0.501}$. Consider the cycle lengths of the permutation $\tau_c$. We claim that there exists a prime number $p < n^{0.501}$ which does not divide any of these lengths. Indeed, otherwise the sum of cycle lengths of $\tau_c$ is not less than the sum of all prime numbers not exceeding $n^{0.501}$ (because each positive integer is not less than the sum of its prime divisors, that in turn follows from $ab \geq a + b$ for $a, b \geq 2$). But for large $n$, the number of prime numbers less than $n^{0.501}$ is at least $n^{0.5009}$ by some weak version of the Prime Number Theorem (that also has many short proofs). Therefore, their sum exceeds $n$, contradiction.

Now we may consider the permutation $\tau_0 \in S_n$ which acts as $\tau_c$ on $C(\tau)$ and has a cycle of length $p$ on $[n] \setminus C(\tau)$. The order of $\tau_0$ is not less than $p \cdot N$, therefore $N \leq f(n)/p \leq f(n)/2$, and by the above argument we get $\mathrm{ord}(\tau) \leq N + R \leq f(n)/2 + n < f(n)$ for large $n$.

\end{document}
